% EagleEye FYP — Daily change report (compiled from git history + session notes)
% Compile: pdflatex EagleEye_Daily_Change_Report.tex (run twice for stable references)
\documentclass[11pt,a4paper]{article}

\usepackage[margin=1in]{geometry}
\usepackage[T1]{fontenc}
\usepackage{lmodern}
\usepackage{microtype}
\usepackage{graphicx}
\usepackage{booktabs}
\usepackage{longtable}
\usepackage{hyperref}
\usepackage{enumitem}
\usepackage{xcolor}
\hypersetup{
  colorlinks=true,
  linkcolor=blue!50!black,
  urlcolor=blue!50!black,
  pdftitle={EagleEye Daily Change Report},
}

\title{\textbf{EagleEye: Intelligent Edge AI Surveillance}\\[0.35em]
\large Daily Change Report and Development History}
\author{Project repository: \texttt{fyp-eagle-eye}}
\date{Report revised: 19 May 2026 \\ \small Synced through the 16 May 2026 model validation work on \texttt{main}}

\begin{document}
\maketitle

\begin{abstract}
This document records project evolution in a professional, date-ordered format. Entries follow \texttt{git log} on \texttt{main}. The May 2026 additions now cover the model-training layout, hard-negative collection, knowledge distillation experiments, and the lessons learned while moving from lab training to ESP32-CAM validation.
\end{abstract}

\tableofcontents
\newpage

\section{Purpose and conventions}
\begin{itemize}[nosep]
  \item \textbf{Scope:} EagleEye spans ESP32-CAM firmware, a Python MQTT--cloud bridge, a React Native (Expo) mobile client, model training assets, and documentation.
  \item \textbf{Sources:} Commit subjects and bodies are paraphrased for clarity; hashes identify the corresponding revision where applicable.
  \item \textbf{Updates:} Append new dated subsections when material changes land in \texttt{main}. Recompile this \texttt{.tex} file to refresh the PDF.
\end{itemize}

%------------------------------------------------------------------
\section{Chronological development summary}

\subsection{25 October 2025}
\textbf{Milestone: repository bootstrap.} Initial README and project upload; early ESP32-CAM bring-up on university Wi-Fi infrastructure.

\subsection{21 November 2025}
\textbf{Milestone: on-device inference.} TensorFlow Lite pipeline validated; interactive verification via DumbDisplay (tap-to-classify prototype).

\subsection{26 November 2025}
\textbf{Maintenance.} Minor commit (no extended message in history).

\subsection{8 December 2025}
\textbf{End-to-end IoT path.} Human detection with approximately 712~ms latency; intruder JPEG relayed via MQTT; Cloudinary storage with metadata URLs persisted in Firebase Realtime Database. Modular layout: camera sketch with \texttt{EagleEye\_IoT.h} for connectivity. \\
\textbf{Security hygiene.} \texttt{.gitignore} extended to exclude secrets (\texttt{.env}, \texttt{secrets.h}, \texttt{serviceAccountKey.json}).

\subsection{10 December 2025}
\textbf{Documentation.} README expanded to describe system architecture and usage.

\subsection{14 December 2025}
\textbf{Local MQTT and payload efficiency.} Mosquitto-based local broker; bridge subscribed for diagnostics; ESP32 transmits raw binary JPEG (larger MQTT buffer, tuned thresholds and cooldown). Supporting assets: Mosquitto configuration, training notebooks, quality simulation utilities, dataset maintenance.

\subsection{10 February 2026}
\textbf{Maintenance.} Minor commit.

\subsection{14 February 2026}
\textbf{Custom tiny model family.} In-house training artifacts under a dedicated model directory; reference to a $\sim$106~ms variant and Arduino export layout.

\subsection{16 February 2026}
\textbf{Greyscale TinyML deployment (\texttt{46847c76}).} INT8 CNN at $48{\times}48{\times}1$; RGB565 capture with software greyscale for inference; color evidence uploads unchanged; tensor arena reduced to 100~kB; simplified capture path in \texttt{EagleEye\_IoT.h}.

\subsection{17 February 2026}
\textbf{Command and control (\texttt{f803947c}).} Mobile: tabbed navigation, dashboard with arm/disarm and latest alert preview, heartbeat-driven online/offline badge. Bridge: Firebase listener on \texttt{config/armed}, intrusion gating when disarmed, periodic \texttt{status/heartbeat} updates. Manual start guide updated for USB/ADB workflow.

\subsection{18 February 2026}
\textbf{Imaging pipeline (\texttt{c79c10f7}, \texttt{719dae9b}).} QVGA capture with center $240{\times}240$ crop prior to model resize---undistorted ``digital zoom'' for improved range; serial diagnostics updated.

\subsection{6 March 2026}
\textbf{Academic milestone.} Mid FYP presentation completed (status commit).

\subsection{21 April 2026}
\textbf{Repository restructure (\texttt{179ef260}).} Consolidated layout: \texttt{firmware/eagleeye-main}, \texttt{backend}, \texttt{mobile-app}, \texttt{model-training}, \texttt{datasets}, \texttt{docs}, \texttt{\_archive}; sketch renamed to \texttt{eagleeye-main.ino}; manuals and README paths corrected. \\
\textbf{Power-aware sensing (\texttt{d87859da}).} PIR on GPIO~14, deep sleep, timed AI windows, armed-state respect, camera/Wi-Fi de-init before sleep; SD disabled to free the PIR pin.

\subsection{2 May 2026 (\texttt{7dc9cc93})}
\textbf{Live preview (first pass).} Git message is short; the merge introduced HTTP streaming from the camera (\texttt{esp\_http\_server}, MJPEG at \texttt{/} on port~80), integration in \texttt{eagleeye-main.ino}, initial \texttt{LiveMonitorScreen.js} (WebView-based), dashboard navigation to Live, \texttt{pir\_debug} test sketch, minor \texttt{bridge.py} edits, and test images under \texttt{backend/captures/}.

\subsection{3 May 2026 (\texttt{3f6989fd})}
\textbf{Live preview reliability + gallery bulk delete.} Firmware: \texttt{eagleeye\_grab\_jpeg()} and \texttt{GET /capture}; WebSocket binary JPEG stream on port~81 (\texttt{eagleeye\_ws.h}, WebSockets library); latch reset when preview ends so AI alerts work again. Mobile: WebSocket + \texttt{expo-image} live path, Android \texttt{usesCleartextTraffic}, remove unused \texttt{react-native-webview}; \textbf{Delete all} alerts with Cloudinary queue via \texttt{deletion\_requests}. Documentation: this \texttt{docs/daily\_changelog/} tree.

\subsection{6 May 2026 (\texttt{0c41c75e})}
\textbf{Model-training hygiene (\texttt{refactor(model-training)}).} Experimental and superseded Python training scripts were moved into \texttt{model-training/\_archive/} (e.g.\ COCO/Colab helpers, colour-model variants, transfer-learning and verification utilities) so the active tree is not cluttered. The single \textbf{production} training entry point for the deployed greyscale TinyML model is consolidated under \texttt{model-training/worlable\_model/train\_tiny\_model.py} (48$\times$48 greyscale, ESP32-oriented). A new \texttt{sketchboard/} tree provides a \textbf{scratch workspace}: copies of the working trainer, a place for firmware experiments, \texttt{notes/}, and a README describing safe day-to-day experimentation without editing production paths.

\subsection{8--9 May 2026}
\textbf{ESP32-CAM Hard Negative Capturer.} Developed a local Python Flask server (\texttt{firmware/tools/hard\_negative\_capturer/}) paired with a custom ESP32 sketch to manually capture false positives/negatives. Images are POSTed over Wi-Fi and sorted into \texttt{Humans}/\texttt{NonHuman} datasets. Addressed system reboots (EV-VSYNC-OVF) and Watchdog Timer timeouts during inference.

\subsection{9--10 May 2026}
\textbf{ML Architecture Evaluation \& Presentation Prep.} Prepared to defend the Kaggle "Overthinkers" submission for CodeRush '26. Focused on validating the multi-model bagging/boosting pipeline, Optuna hyperparameter optimization, and adaptive blending logic.

\subsection{11 May 2026}
\textbf{Hard Negative Retraining.} Fine-tuned the custom tiny model in \texttt{sketchboard/} to integrate hard negative datasets, targeting 90\%+ accuracy. Added data augmentation (flip, rotate, zoom) in \texttt{train\_in\_sketchboard.py}. Exported updated INT8 \texttt{.tflite} models and generated C-headers (\texttt{human\_detect\_model\_data\_augmented.h}) for firmware integration.

\subsection{16 May 2026}
\textbf{Distillation and validation pass.} The work moved beyond the earlier single hard-negative retraining loop and tested a more systematic model-improvement path. Fresh ESP32-CAM validation captures were added under \texttt{sketchboard/dataset/}, the exported model header was refreshed in \texttt{sketchboard/human\_detect\_model\_data.h}, and the training/evaluation scripts under \texttt{sketchboard/knowledge\_distillation/} were used to compare a larger teacher model with a smaller ESP32-friendly student model.

\textbf{Changed files.} Key files for this day include \texttt{sketchboard/dataset/*20260516*.jpg}, \texttt{sketchboard/knowledge\_distillation/train\_teacher.py}, \texttt{train\_distillation.py}, \texttt{train\_kd\_optuna.py}, \texttt{evaluate\_tflite.py}, \texttt{sketchboard/train\_optuna\_cv.py}, \texttt{sketchboard/sketchboard\_hard\_negative\_90plus\_metrics.json}, and the generated model header used by sketchboard firmware tests.

\textbf{New method.} Instead of only adding more hard-negative images and retraining the same tiny network, this pass used a teacher/student knowledge distillation workflow. Optuna was introduced to search for better distillation settings such as learning rate, dropout, temperature, and alpha, making the model selection process more explainable than manual tuning.

\textbf{Result / lesson.} The experiment showed a practical direction for improving accuracy without forgetting the ESP32-CAM constraint: the student model still needs to stay small and fast, but it can learn from a stronger teacher and be judged with repeatable metrics rather than only live ad-hoc tests.


%------------------------------------------------------------------
\section{Suggested maintenance}
\begin{itemize}[nosep]
  \item After each merged feature, append a dated subsection (one paragraph + bullet list if needed) and re-run \texttt{pdflatex}.
  \item Keep commit messages descriptive; this report is easier to maintain when subjects follow \texttt{type(scope): summary} convention.
\end{itemize}

\section{Compilation}
\texttt{cd docs/daily\_changelog} \\
\texttt{pdflatex EagleEye\_Daily\_Change\_Report.tex} \\
\texttt{pdflatex EagleEye\_Daily\_Change\_Report.tex}

\end{document}
